% TOP_PROBLEM: {"id": 3388, "title": "A sextic counterexample to Euler's sum of powers conjecture", "category_id": 1, "category": {"id": 1, "name": "number_theory", "display_name": "Number Theory", "description": "Properties of integers, prime numbers, Diophantine equations.", "slug": "number-theory", "order_index": 1, "created_at": "2026-07-31T15:26:25.670Z"}, "status": "open", "problem_number": "OPG-508", "set_id": 12}
% FIRST_POSTED: 2026-10-01
% ENTRY_KIND: statement_only
% UnsolvedMath ID 3388; source catalogue code OPG-508
% !TeX program = lualatex
% Definition and sources only; this file is not an LLM solution attempt.
\documentclass[11pt,a4paper]{article}
\usepackage[margin=25mm]{geometry}
\usepackage{fontspec}
\setmainfont{lmroman10-regular.otf}[BoldFont=lmroman10-bold.otf,
 ItalicFont=lmroman10-italic.otf,BoldItalicFont=lmroman10-bolditalic.otf]
\usepackage{amsmath,amssymb,mathtools}
\usepackage{unicode-math}
\setmathfont{latinmodern-math.otf}
\AtBeginDocument{\let\setminus\smallsetminus}
\usepackage{xurl,hyperref}
\hypersetup{colorlinks=true,urlcolor=blue}
\setcounter{secnumdepth}{0}
\setlength{\parindent}{0pt}
\setlength{\parskip}{5pt}
\setlength{\emergencystretch}{3em}
\begin{document}
\section*{A sextic counterexample to Euler's sum of powers conjecture}\label{3388}
{\small UnsolvedMath ID 3388 · OPG-508 · Number Theory · OpenGarden}
\subsection{Definitions and mathematical statement}
The Diophantine existence problem is to find positive integers
$x_1,\ldots,x_6$ satisfying
\[
             x_1^6+x_2^6+x_3^6+x_4^6+x_5^6=x_6^6,
\]
or to prove that there is no such tuple. Every variable is strictly
positive, so no zero term can be used to reduce the number of summands.
The five integers on the left need not be distinct. Permuting them
does not change the equation, and one may impose
$x_1\le\cdots\le x_5<x_6$ when searching.

Because the equation is homogeneous, multiplying all entries by the
same positive integer preserves it. Existence is therefore equivalent
to existence of a primitive tuple, meaning
$\gcd(x_1,\ldots,x_6)=1$; divide by the common divisor to obtain one.
It is also equivalent to the existence of positive rational numbers
$y_1,\ldots,y_5$ with $\sum_i y_i^6=1$, by normalizing $x_6$ and
then clearing denominators in the reverse direction.
The question fixes both the exponent six and the number five of
summands. Examples with a different exponent, an extra summand,
or some zero coordinates do not satisfy this target. A positive
answer can be certified by direct exact integer arithmetic.
\subsection{Short English statement}
Can a positive sixth power equal the sum of exactly five positive sixth powers, with repeated summands allowed?
\subsection{Sources}
\begin{itemize}
\item[S1] ULAM.AI and the UnsolvedMath Contributors, supplied problems.json, record 3388 (OPG-508), OpenGarden collection. Catalogue identity and source transcription; CC BY 4.0.
\url{https://huggingface.co/datasets/ulamai/UnsolvedMath}.
\item[S2] Open Problem Garden, A sextic counterexample to Euler's sum of powers conjecture. Original problem and discussion; the mathematical formulation below is expanded from the supplied source record.
\url{http://www.openproblemgarden.org/op/a_sextic_counterexample_to_eulers_sum_of_powers_conjecture}.
\end{itemize}
\end{document}
